\chapter{Konforme Verfeinerung und Fehlerkontrolle}\label{theory:adaptive}
Ein affines Tetraeder wird an einer Kante \((a,b)\) durch den Mittelpunkt
\(m=(a+b)/2\) in zwei Tetraeder geteilt. Je ein Kind ersetzt den Eckpunkt
\(a\) beziehungsweise \(b\) durch \(m\). Beide Kinder liegen im Elternsimplex;
ihre Volumina summieren sich zum Elternvolumen. Dieselbe Konstruktion gilt für
Randdreiecke und ihre Flächen. Jeder vollständige Kantenstern erhält denselben
Mittelpunkt. Auf einer gemeinsamen Fläche stimmen daher Teilkanten und Teilflächen
überein: Es entstehen keine hängenden Knoten. Mehrere ausgewählte Kanten werden
in gemeinsamer deterministischer Reihenfolge verarbeitet.

Markierung verwendet das vorhandene Residualmaß \(\eta_K\) und die minimale
Dörfler-Menge, deren quadrierte Summe mindestens \(\theta\sum_K\eta_K^2\) erreicht.
Gleichstände werden nach Eingabezellzeilen entschieden. Die längste Kante jeder
markierten Eingabezelle wird geteilt; der notwendige Sternabschluss kann auch
unmarkierte Zellen betreffen. Dieser Vertrag liefert weder garantierte globale
Formregularität noch zertifizierte Zuverlässigkeits-/Effizienzkonstanten des
Indikators. Degeneration, verlorene Gruppen und ungültige Elterngeometrie sind
Abbruchfehler. Die tatsächliche Genauigkeit muss durch unabhängige Referenzfälle
und darf nicht allein durch abnehmendes \(\eta\) beurteilt werden.
